\documentclass[12pt]{article}
\usepackage[utf8]{inputenc}
\usepackage[T1]{fontenc}
\usepackage{hyperref}
\title{Nobiin Language Resources for K-12 Curriculum Development}
\author{Compiled by Claude Sonnet 5.5 (claude-sonnet-5-5) for Andr\'as Kornai}
\date{}
\begin{document}
\maketitle

\section{Introduction}
Nobiin (ISO 639-3: fia) is a Nubian language of Sudan (and Egypt), tonal, with contrastive vowel and consonant length and SOV order. It has no standardized orthography; it has been written in Latin and Arabic scripts, and there are recent efforts to revive the Old Nubian alphabet. This search was brief and found few resources.

\section{Primary Educational Resources}
\begin{itemize}
\item ``Learn Nubian! (Nobiin) Premium'', a paid iPhone app with about 20 lessons of vocabulary and grammar and script tracing: \url{https://apps.apple.com/ph/app/learn-nubian-nobiin-premium/id6476841617}
\end{itemize}
No printed primer or textbook was found in the results.

\section{Linguistic and Reference Materials}
\begin{itemize}
\item Wikipedia, Nobiin language; it states that the Latin orthography follows Roland Werner's 1987 \emph{Grammatik des Nobiin}, described there as the only published grammar (a claim that may be dated): \url{https://en.wikipedia.org/wiki/Nobiin_language}
\end{itemize}
No Nobiin dictionary was found in the results.

\section{Religious and Community Resources}
None found in the search results.

\section{Digital and Audio Resources}
Only the app listed above.

\section{Curriculum Development Context}
The absence of a standard orthography is the central obstacle; an orthography decision (Latin, Arabic or Old Nubian-based) would precede materials development. Werner's grammar is the main lead for a linguistic base. Further searching (library catalogs, Sudanese and Egyptian Nubian cultural organizations) is recommended.
\end{document}